\subsubsection{Separable vs Entangled States}

Let the general state of a two-qubit system be (as defined in \autoref{eq:general-two-qubit-state}, \autopageref{eq:general-two-qubit-state}):
\begin{equation*}
    \ket{v} = c_{0}\ket{00} + c_{1}\ket{01} + c_{2}\ket{10} + c_{3}\ket{11}
\end{equation*}
The question is: \hl{\emph{can we interpret this as two independent qubits?}} The answer is \textbf{not always}. As we saw on \autopageref{subsubsec:tensor-product}, the answer depends on whether the states are separable or entangled.

\highspace
\begin{flushleft}
    \textcolor{Green3}{\faIcon{book} \textbf{Separable States}}\index{When is a state separable?}
\end{flushleft}
A state is \textbf{separable if it can be written as a tensor product of two single-qubit states}. In other words, there exist single-qubit states $\ket{v_{A}}$ and $\ket{v_{B}}$ such that:
\begin{equation*}
    \ket{v} = \ket{v_{A}} \otimes \ket{v_{B}}
\end{equation*}
So the system can be split into one state for qubit $A$ and one state for qubit $B$, and the global state is just the combination of these two states. As consequence, the coefficients factorize as:
\begin{equation*}
    c_{0} = a_{0}b_{0}, \quad c_{1} = a_{0}b_{1}, \quad c_{2} = a_{1}b_{0}, \quad c_{3} = a_{1}b_{1}
\end{equation*}

\highspace
\begin{flushleft}
    \textcolor{Green3}{\faIcon{book} \textbf{Entangled States}}\index{When is a state entangled?}
\end{flushleft}
A state is \textbf{entangled if it cannot be written as a tensor product of two single-qubit states}. In other words, there do not exist single-qubit states $\ket{v_{A}}$ and $\ket{v_{B}}$ such that:
\begin{equation*}
    \ket{v} \neq \ket{v_{A}} \otimes \ket{v_{B}}
\end{equation*}
In this case, the state cannot be separated into two independent qubits, and the coefficients do not factorize as above. \hl{The system is simply one single entity.} For example, the Bell state (we will see in future sections) is an entangled state:
\begin{equation*}
    \ket{\Phi^{+}} = \frac{1}{\sqrt{2}}\left(\ket{00} + \ket{11}\right)
\end{equation*}
And the demonstration was written on \autopageref{ex:non-separable}. In other words, we can only say that the system is in state $\ket{\Phi^{+}}$, but we cannot say that qubit $A$ is in some state and qubit $B$ is in some state, because the state of the system does not factorize into two independent states.

\highspace
At this stage, entanglement is introduced as a structural property of quantum states. Its physical meaning and implications become clear only when studying measurement and quantum circuits in later sections.

\highspace
\textcolor{Green3}{\faIcon{balance-scale} \textbf{Difference.}} With \textbf{separable} states, the \hl{information is stored in the individual qubits}. In contrast, with \textbf{entangled} states, the \hl{information is stored in the correlations between the qubits}.